% !TeX root = ../papers/nc-p3.tex
\section{一般の旗を用いた反対階数10の有限片側対称枝の完結}
\label{sec:finite-rank-ten-full-regularity}
\begin{theorem}\label{thm:finite-rank-ten-full-regularity}
標準有限Frobenius窓の全辺が有限で、一方の反対交換が対称かつ反対階数が10とする。
左右Koszul AS正則性の必要十分条件は、二つの奇位置のインスタントンモナドがともに束を与えることである。
従って定理\ref{thm:finite-elliptic-octic-windows}の開条件を満たす全データが正則である。
四次Grassmann写像による添字保持同型・周期の判定はこの枝の特殊値にも適用できる。
\end{theorem}
\begin{proof}
必要性は前の境界排除で得た。十分性では窓を一つ固定する。
以下で一般に選ぶのはこの固定された写像の旗であり、窓そのものを一般化しない。

\paragraph{次数二の対応の全点での定義。}
片側の一般切断の零曲線 $F$ を固定し、$L=\mathcal O_F(h)$ とする。
補題\ref{lem:octic-linear-determinant}の $J\in\mathrm{Pic}^4(F)$ に対する四次行列式を $q_J$ と書く。
これらは $\mathbb P H^0(I_F(4))=\mathbb P^2$ の中の直線をちょうど二重に覆い、恒等的に零になる $q_J$ はない。
実際、幾何学的なPicard曲線上で正規化した普遍線束を使い、
$H^0(J)$ と $H^0(LJ)$ の束を $B_4,B_{12}$ と書く。階数は添字通りで、次数はともに $-1$。
後者の次数公式は楕円曲線の積の正規化Poincaré束の第一Chern類の平方の繊維積分が
点類の $-2$ 倍であることからRiemann--Rochで従う
（第一Chern類は対角から座標二繊維を引いた類である）。
従って $K_J$ の束の次数は $4(-1)-(-1)=-3$ で、$q_J(x)$ は $J$ の関数として
固定した次数2の線束の切断となる。もし係数が一次元しか張らなければ行列式は
$q(x)e(J)$ に分かれる。非零四次式 $q$ の零曲面上で全 $J$ の評価が失敗するが、
同補題の証明では、このような $x\in\mathbb P^3$ は $F$ と高々一次元の平面射影例外集合に限られた。
従って係数は二次元全部を張り、次数2の完全鉛筆は基点を持たない。
また $x\in F$ で評価した関係は $J_x$ で零になる三次元部分に入るので $q_J$ は $F$ を含む。
これで直線とその二重被覆を得る。

$F$ を与える切断を含む二平面を動かすと対応する $T$ は上の網の全てを動く。
実際、切断と第二切断の行列式を取る写像はSerreの評価列より
切断の四空間を第一切断で割った空間から $H^0(I_F(4))$ への同型である。
一般の旗では $T$ は行列式の直線の外にある。よって全 $J$ で $q_J|_T$ は恒等的に零ではなく、
その $F$ を除いた因子は $|C|$ の一繊維である。
得られた写像 $g_F:\mathrm{Pic}^{-4}(F)\to |C|$ は逆線束を取ってから上の直線を $[T]$ から射影したもので次数2。
ここで残余因子が定理\ref{thm:octic-fiber-twist-transfer}の余核の行列式に等しいことは、
$W\subset H^0(J)\otimes\mathcal O_T$ の行列式が $\mathcal O_T(-F)$ であることからも従う。
二側を交換してもこの旗の開条件は成立する。よって一般の $T$ とその両鉛筆の一般の $F,C$ で同時に使える。

\paragraph{二つの繊維制限。}
$U$ を滑らかな楕円繊維とそのPicard曲線の普遍層による変換とする。
始域にねじれを入れれば基点を選ばず定義でき、幾何学的には通常の楕円曲線の同値である
（前節の相対対角イデアルの構成を一繊維に適用してもよい）。
$F$ で次数 $-4$、$C$ で次数12を使い、
\[
 I_F=i_{F*}U,\quad r_F=U^{-1}\mathbf L i_F^*,\qquad
 I_C=A_0 i_{C*}U,\quad r_C=U^{-1}\mathbf L i_C^* A_0^{-1}
\]
とする。$r$ はそれぞれ $I$ の左随伴である。
$D_F=L_F,\ D_C=A_0 L_C A_0^{-1}$ と書けば $\Psi=\Phi^{-1}[2]=D_F D_C$ である。
$D_F$ は $r_F$ 上で恒等、$r_C$ 上で次数2の線束のテンソル積として働く。
$D_C$ は $r_C$ 上で恒等、$r_F$ 上で次数2の線束のテンソル積として働く。

この最後の主張は各点の数値作用だけからは従わないので族で確認する。
$A_0^{-1}$ を $I_F$ の普遍核に適用すると、前段により全パラメータで
同じ移動の行列式余核になる。各行列は束間の単射で、余核はパラメータに平坦である。
随伴行列によりその余核は $T\times_{|C|}\mathrm{Pic}^{-4}(F)$ にスキームとして支えられる。
従って $L_C$ はこの核に $g_F^*\mathcal O(1)$ を掛けるのと同じであり、
$D_C I_F=I_F(-\otimes g_F^*\mathcal O(1))$、左随伴により所望の $r_F D_C$ の式となる。
逆方向には命題\ref{prop:quartic-k3-two-pencil-twists}の $B_0A_0=L_{-4h}[2]$ より
$A_0=L_{-2h}A_0^{-1}L_{-2h}[2]$ であることを使う。次数12の線束を $-2h$ 捩れば次数 $-4$ なので、
$C,F$ を交換した同じ行列式で、$A_0 I$ の相手繊維への台の写像も次数2となる。
同じ側の作用は繊維の法線束が自明なことから従う。
この証明は一般の印のPicard曲線にわたる全点の族を使い、相手の特異繊維での安定性を必要としない。

\paragraph{制限の階数、次数、半安定性。}
前節の $P_i$ を用いる。$r_F P_i$ は同じ移動を持つ階数4の束である。
$i=0,1$ の全 $J$ でのコホモロジーは前節と補題\ref{lem:octic-linear-determinant}の通りである。
$i=2,3$ には
\[
 \Psi P_2=M_h^*(D+h)[1],\qquad \Psi P_3=\mathcal O(D+h)[1]
\]
を使う。実際 $T_{G_2}^{-1}G_2=G_2[1]$ と双対評価列を左変異表示に代入すればよい。
$\Psi$ は既に $r_F$ 上で線束の捩りなので、右辺で共通の移動を検査すれば足りる。
$D+h$ の $F$ 制限は $-h$ であり、次数 $-4$ の線束に $h$ を足すと次数4になる。
その $M_h$ との積の $H^1$ は同補題で零、$H^0$ は四次元である。
従って四対象とも $\mathbf R\mathrm{Hom}(P_i,I_F k(z))$ は全点で次数1に四次元だけを持つ。
一方 $A_0^{-1}P_i$ は順に
$M_h^*(h),\mathcal O(h),\mathcal O(D)[-1],M_h(D)[-1]$。
$D|_C=2h|_C$ なので、次数12の線束に対して同じ計算が全点で次数零の四次元を与える。
コホモロジーと基底変換により $r_C P_i$ も共通の移動の階数4束となる。

各側の一般の旗の関数体 $K$ で、印を付けた楕円繊維の全ての $K$ 線束の次数は8の倍数である。
実際、旗は切断の直線 $s$ とそれを含む二平面 $P$ からなる。
普遍零繊維の全空間は $x\in\mathbb P^3$、$s\subset f(x)$、$s\subset P\subset H$ を順に選ぶ射影束塔である。
従ってPicard群は $h$ と二つの相対同語反復線束で生成され、後二者は $s,P/s$ による旗底からの引戻しである。
滑らかな一般繊維の線束は有理切断の因子の閉包をこの正則全空間で取って延長でき、繊維に残るのは $h$ のみである。
反対側は双対Grassmann写像で同じである。

従って $K$ 上の対応するねじれPicard曲線の任意の層の階数は4の倍数である。
実際 $U$ による像を楕円繊維上で $(r,d)$ と書けば $d$ は8の倍数で、幾何学的点層とのEuler対合は
$-4r-d$ または $12r-d$ である。この対合がねじれ層の階数である。
各束の幾何学的最大不安定部分は一意であり、エタール局所の表示で降下して $K$ 上のねじれ部分層になる。
階数4未満の正階数を排除したので四束とも幾何学的半安定である。
この性質は一般の旗へ開に広がる。幾何学的Picard曲線で普遍線束を正規化したとき、
次数 $e$ の変換 $U$ の数値行列は階数・次数順で
$\left(\begin{smallmatrix}0&1\\-1&e\end{smallmatrix}\right)$ である
（点は次数 $e$ の線束、構造層は符号反転した長さ1の点類へ移る）。
従って共通移動を除いた $r_F P_i$ の次数は $(-1,-3,-5,-7)$、$r_C P_i$ の次数は $(3,1,-1,-3)$ である。
実際 $P_i|_F$ の階数は $(1,3,5,7)$、次数は $(0,-8,-16,-24)$、
$(A_0^{-1}P_i)|_C$ の階数は $(3,1,-1,-3)$、次数は $(32,8,-16,-40)$ である。
自明化の変更は次数に共通の4の倍数を加えるだけである。
奇数次数なのでこれらは安定であり、両側で傾きの差は $(i-j)/2$ である。

\paragraph{全次数への伝播。}
$Z_m^{ij}=\mathbf R\mathrm{Hom}_T(P_i,\Psi^mP_j)$ とする。
$Q\to D_C Q\to I_C r_C Q$ と $Q\to D_F Q\to I_F r_F Q$ は
それぞれ一般の $C,F$ の繊維方程式による三角である。
$P_i$ から射を取った第三項は楕円Picard曲線上の $r_C$ または $r_F$ の間の導来Homそのものである。
従って $Z_m^{ij}$ から $Z_{m+1}^{ij}$ への二段で加わる項は、
同じ移動の二束間の導来Homであり、その傾き差は順に
\[
 (i-j)/2+2m,\qquad (i-j)/2+2m+2.
\]
ここで先に $D_C$ を掛ける。二関手の可換性は使っていない。
$Z_0^{ij}$ は前節の随伴三角から負次数と次数1で消える。
従って二段を経ても負次数のコホモロジーは出ない。
最初の周で $i<j$ のときは前節の内部次数 $4+i-j\le3$ の議論と随伴三角から $Z_1^{ij}$ の正次数も消える。
$i>j$ のときは初期の次数2も消え、二段とも正の傾き差なので楕円曲線上の $H^1$ は消える。
$i=j$ では第一段の連結写像は、Serre双対の下で恒等射を安定な制限束の恒等射へ送る制限の双対であり同型である。
これは $A_0^{-1}$ で繊維制限の三角へ戻し、繊維の法線束の自明性を使えばよい。
よって次数1と2が消え、次の段では正の傾き差がある。
以後 $m\ge1$ ならいずれの傾き差も正である。よって全 $m\ge1$ で $Z_m^{ij}$ は純粋となる。
前節の $X_m^{ij}\to Z_m^{ij}\to X_{m-1}^{ij}[1]$ は随伴によるので格子の仮定なしに使え、
帰納的に全Serre純粋性を得る。これで十分性と枝の判定が従う。
\end{proof}
